% Generated draft: requires mathematical and source review.
\documentclass[11pt]{article}
\usepackage[margin=25mm]{geometry}
\usepackage{fontspec}
\setmainfont{Latin Modern Roman}
\usepackage{amsmath,amssymb,mathtools,unicode-math}
\setmathfont{Latin Modern Math}
\usepackage{enumitem,needspace,xurl,hyperref}
\setlength{\parindent}{0pt}
\setlength{\parskip}{5pt}
\setlength{\emergencystretch}{3em}
\newcommand{\N}{\mathbb{N}}
\newcommand{\Z}{\mathbb{Z}}
\newcommand{\Q}{\mathbb{Q}}
\newcommand{\R}{\mathbb{R}}
\newcommand{\C}{\mathbb{C}}
\newcommand{\F}{\mathbb{F}}
\newcommand{\A}{\mathbb{A}}
\newcommand{\PP}{\mathbb{P}}
\newcommand{\E}{\mathbb{E}}
\newcommand{\eps}{\varepsilon}
\newcommand{\dd}{\,\mathrm{d}}
\begin{document}
\section*{UnsolvedMath ID: 30002194}
\textbf{Source problem code:} catalogue entry
\subsection{Definitions and mathematical statement}
The supplied record groups six research questions under one problem. Several are requests for a classification or an explanation, and the record does not supply enough hypotheses to turn all six into uniquely determined mathematical propositions.

For terminology, let \(k\) be a field and let \((V,c)\) be a braided vector space over \(k\). Its Nichols algebra is
\[
\mathcal B(V)=T(V)/I(V),
\]
where \(T(V)\) is the braided tensor Hopf algebra with \(V\) primitive and \(I(V)\) is the largest homogeneous braided Hopf ideal contained in \(\bigoplus_{n\geq 2}T^n(V)\). The braiding is of diagonal type if there is a basis \(x_1,\ldots,x_\theta\) of \(V\) and scalars \(q_{ij}\in k^\times\) such that
\[
c(x_i\otimes x_j)=q_{ij}x_j\otimes x_i.
\]
Finite dimensionality of the Nichols algebra means \(\dim_k\mathcal B(V)<\infty\). These definitions do not specify a characteristic or a ground field for the questions below; the supplied record leaves those choices unresolved.

A groupoid is a category in which every morphism is invertible. Weyl groupoids in the questions are groupoids associated with systems of root data and generated by their simple reflections. The supplied record does not specify the admissible root data, the meaning of finiteness, or the equivalence relation to be used in a classification. Consequently, the term ``finite Weyl groupoid'' is retained with its source-dependent scope unresolved.

For context only, in the real central convention a finite hyperplane arrangement in a real vector space \(E\) is a finite set \(\mathcal A\) of linear hyperplanes. Its chambers are the connected components of
\[
E\setminus\bigcup_{H\in\mathcal A}H.
\]
An essential arrangement is simplicial if each chamber is an open simplicial cone. The record does not establish that this real, central, essential convention is the intended scope of its arrangement question.

The six recorded questions are as follows.
\begin{enumerate}
\item \textbf{Realization by Nichols algebras.} Does every finite Weyl groupoid occur as a Weyl-groupoid invariant of some Nichols algebra? Schematically, the proposed universal realization statement is
\[
\text{for every finite Weyl groupoid }\mathcal W,\quad
\text{there exists a Nichols algebra }\mathcal B(V)
\text{ with }\mathcal W(\mathcal B(V))\simeq\mathcal W.
\]
Here \(\mathcal W(\mathcal B(V))\) denotes the intended associated invariant only when it is defined. The admissible Nichols algebras, the construction of this invariant, and the meaning of \(\simeq\) are not specified. In particular, the source does not restrict this question to finite-dimensional Nichols algebras or to diagonal type.
\item \textbf{Simplicial arrangements.} Classify simplicial arrangements. The ambient field, dimension, central or affine setting, additional restrictions, and classification equivalence are not supplied.
\item \textbf{Cluster algebras.} Determine how far the connection between reflection groupoids and cluster algebras extends. Neither the existing connection to be extended nor the desired correspondence or theorem is specified. The record uses ``reflection groupoids'' here; it does not explicitly identify their scope with that of the Weyl groupoids in the other questions.
\item \textbf{Toric varieties.} Clarify the connection with toric varieties. The source does not specify the objects to be related, a construction, a class of toric varieties, or a proposed equivalence.
\item \textbf{Sporadic cases.} Explain the sporadic finite Weyl groupoids. The record supplies neither a list of sporadic cases nor a mathematical criterion for what would constitute an explanation.
\item \textbf{Derivation of a classification.} Determine whether the classification of finite-dimensional Nichols algebras of diagonal type can be derived from finite Weyl groupoids. The field, characteristic, equivalence relation for Nichols algebras, and intended meaning of ``derived'' remain unspecified. This is a request for a classification method or deduction, rather than a stated implication with specified premises and conclusion.
\end{enumerate}
These questions are retained together as the supplied research program. The unresolved qualifications above prevent a single fully specified theorem or conjecture from being extracted without additional source information.
\subsection{Short English statement}
The record asks six related questions about realizing finite Weyl groupoids by Nichols algebras, classifying simplicial arrangements, extending connections with cluster algebras and toric varieties, explaining sporadic cases, and deriving a classification of finite-dimensional Nichols algebras of diagonal type. Several questions lack the hypotheses or target correspondence needed for a precise proposition. The supplied ``partially solved'' label is dated curation, not an independently verified status.
\subsection{Sources}
\begin{itemize}
\item Supplied local UnsolvedMath registry record, permanent ID \(30002194\), exact source problem code \texttt{catalogue entry}, title ``Classification Problems for Finite Weyl Groupoids,'' in the collection ``Oberwolfach Reports Open Problems.'' No individual proposer is supplied; the proposed year is recorded as 2012.
\item Supplied report citation: \emph{Mini-Workshop: Nichols Algebras and Weyl Groupoids (2013)}, Oberwolfach Reports, DOI 10.4171/owr/2012/47, pp.~2879--2905. \url{https://doi.org/10.4171/owr/2012/47}. The record attributes these questions to workshop open-problem discussions and states that the official report lists all six together. Its \texttt{corrected\_verified} statement label and verification assertion are supplied registry metadata, not independent verification here. The citation's 2013 designation and the record's proposed year 2012 are retained without resolving their relationship.
\item Supplied literature source, ``original OWR report'': \url{https://doi.org/10.4171/owr/2012/47}.
\item Supplied literature source, Cuntz--Heckenberger, \emph{Finite Weyl groupoids}: \url{https://arxiv.org/abs/1008.5291}.
\item Supplied literature source, Cuntz--Lentner, \emph{A simplicial complex of Nichols algebras}: \url{https://arxiv.org/abs/1503.08117}.
\item The exact supplied literature-check date is \texttt{2026-08-21}. The supplied dated curation labels the record ``partially solved'' and reports that finite Weyl groupoids and several Nichols-algebra classifications are known, while the six questions include unresolved realization, simplicial-arrangement, cluster, toric, sporadic-structure, and deduction questions. This assessment is supplied dated curation derived from the OWR material; neither its time-sensitive status nor its cited claims have been independently verified here.
\end{itemize}
\end{document}
